\section{Software: Freefem++}
\label{sec:WP7:Freefem++:software}

\begin{table}[h!]
    \centering
    { \setlength{\parindent}{0pt}
    \def\arraystretch{1.25}
    \arrayrulecolor{numpexgray}
    {\fontsize{9}{11}\selectfont
    \begin{tabular}{!{\color{numpexgray}\vrule}p{.4\textwidth}!{\color{numpexgray}\vrule}p{.6\textwidth}!{\color{numpexgray}\vrule}}
        \rowcolor{numpexgray}{\rule{0pt}{2.5ex}\color{white}\bf Field} & {\rule{0pt}{2.5ex}\color{white}\bf Details} \\
        \rowcolor{white}\textbf{Consortium} & \begin{tabular}{l}
Sorbonne U\\
\end{tabular} \\
        \rowcolor{numpexlightergray}\textbf{Exa-MA Partners} & \begin{tabular}{l}
Inria PARIS\\
Sorbonne U\\
\end{tabular} \\
        \rowcolor{white}\textbf{Contact Emails} & \begin{tabular}{l}
frederic.hecht@sorbonne-universite.fr\\
pierre-henri.tournier@sorbonne-universite.fr\\
pierre.jolivet@sorbonne-universite.fr\\
\end{tabular} \\
        \rowcolor{numpexlightergray}\textbf{Supported Architectures} & \begin{tabular}{l}
CPU Only\\
\end{tabular} \\
        \rowcolor{white}\textbf{Repository} & \href{https://github.com/FreeFem/FreeFem-sources}{https://github.com/FreeFem/FreeFem-sources} \\
        \rowcolor{numpexlightergray}\textbf{License} & \begin{tabular}{l}
OSS:: LGPL v*\\
\end{tabular} \\
        \bottomrule
    \end{tabular}
    }}
    \caption{WP7: Freefem++ Information}
\end{table}

\subsection{Software Overview}
\label{sec:WP7:Freefem++:summary}

FreeFEM++ is a mature finite-element prototyping environment widely used for rapid PDE experimentation. In WP7 it provides a
lightweight yet expressive vehicle to validate the benchmarking methodology on CPU-focused applications: scripts can be
containerised, executed through WP7 ReFrame suites, and paired with ParMmg mesh-adaptation services originating from WP1. The
software is also used to cross-check algorithmic developments from WP3 (domain decomposition) before they are ported to large
scale codes such as HPDDM.

\begin{table}[h!]
    \centering
    { 
        \setlength{\parindent}{0pt}
        \def\arraystretch{1.25}
        \arrayrulecolor{numpexgray}
        {
            \fontsize{9}{11}\selectfont
            \begin{tabular}{!{\color{numpexgray}\vrule}p{.25\linewidth}!{\color{numpexgray}\vrule}p{.6885\linewidth}!{\color{numpexgray}\vrule}}
    
    \rowcolor{numpexgray}{\rule{0pt}{2.5ex}\color{white}\bf Features} &  {\rule{0pt}{2.5ex}\color{white}\bf Short Description }\\ 
    
\rowcolor{white}    Geophysics & Seismic wave propagation and subsurface imaging prototypes jointly developed with WP4 inversion teams. \\
\rowcolor{numpexlightergray}    Energy & Thermal-fluid studies for nuclear safety scenarios and heat exchanger optimisation. \\
\rowcolor{white}    Health & Biomedical diffusion and electrophysiology models used for teaching and reduced-order benchmarks. \\
\rowcolor{numpexlightergray}    Environment & Coastal erosion and pollutant transport simulations leveraging adaptive meshing. \\
\rowcolor{white}    Energy & Magnetohydrodynamic stability analyses supporting fusion-device design studies. \\
\rowcolor{numpexlightergray}    Aero & Incompressible and compressible aerodynamics mini-apps (ParMmg-coupled) for mesh adaptation workflows. \\
\rowcolor{white}    mini-apps & Collection of ParMmg/FreeFEM coupled kernels packaged as WP7 demonstrators (mesh smoothing, anisotropic refinement). \\
\end{tabular}
        }
    }
    \caption{WP7: Freefem++ Features}
\end{table}


\subsection{Parallel Capabilities}
\label{sec:WP7:Freefem++:performances}


FreeFEM++ offers an MPI-enabled variant that links against PETSc, HPDDM, and MUMPS for large-scale sparse solves. Within WP7 the
code is exercised on TGCC's IRENE Rome partition and the IDRIS Jean Zay CPU nodes, reaching up to 2048 cores for mesh-adaptation
studies. Thread-level parallelism is leveraged via OpenMP inside core kernels, and experimental HIP back-ends are evaluated for
selected operators. The runtime integrates smoothly with WP7 container images and uses shared Spack recipes to ensure that HPDDM
and ParMmg dependencies match those of the larger Exa-MA ecosystem.


\subsection{Initial Performance Metrics}
\label{sec:WP7:Freefem++:metrics}

WP7 benchmarking of FreeFEM++ centres on the ``app-freefem-parmmg'' workflow coupling anisotropic mesh adaptation with domain
decomposition solvers. Input meshes (3~GB total) and configuration scripts are distributed through the WP7 FreeFEM++ Zenodo workspace (release in preparation).
Runs on IRENE (48-core nodes) show 88\% strong-scaling efficiency between 4 and 64 nodes, while hybrid OpenMP+MPI configurations
reduce assembly time by 22\%. Energy profiles collected via RAPL reveal a 15\% improvement when adopting WP7's NUMA-aware
launcher. Current limitations are tied to sequential preprocessing steps in the mesh adaptation pipeline and to filesystem
contention when writing intermediate meshes; mitigations are being co-designed with WP7's data-management task.

\subsubsection{Benchmark \#1}
\begin{itemize}
    \item \textbf{Description:} ``app-freefem-parmmg'' anisotropic mesh adaptation on IRENE Rome nodes (up to 64 nodes) with $80$~million elements targeting faulted reservoir geometries.
    \item \textbf{Benchmarking Tools Used:} ReFrame orchestrates batches; Extrae profiles MPI communication; WP7 energy plug-in records node power; filesystem counters captured via Darshan.
    \item \textbf{Input/Output Dataset Description:}
        \begin{itemize}
            \item \textbf{Input Data:} Geological mesh pack (3~GB, Gmsh format) hosted on the WP7 FreeFEM++ Zenodo workspace.
            \item \textbf{Output Data:} Refined meshes (Gmsh + VTK) and convergence metrics archived with accompanying JSON metadata.
            \item \textbf{Data Repository:} Zenodo community ``FreeFEM++/ParMmg Exa-MA datasets'' curated by WP1 and WP7 engineers.
        \end{itemize}
    \item \textbf{Results Summary:} 88\% strong-scaling efficiency from 4 to 64 nodes; average energy-to-solution reduced by 15\% via WP7 NUMA pinning policies; mesh quality metrics improved by 2.3× relative to isotropic refinement.
    \item \textbf{Challenges Identified:} Intermediate mesh dumps saturate Lustre metadata servers; ongoing work tests batched checkpoints and DAOS staging.
\end{itemize}

\subsection{12-Month Roadmap}
\label{sec:WP7:Freefem++:roadmap}

\begin{itemize}
    \item \textbf{Data improvements:} Automate mesh-pack publication with provenance metadata and reproducibility badges for each benchmark (Q1~2026).
    \item \textbf{Methodology application:} Integrate FreeFEM++ scripts into WP7 GitHub Actions (container build + ReFrame smoke tests) covering HPDDM-enabled workflows (Q2~2026).
    \item \textbf{Results retention:} Produce a quarterly performance digest combining scaling plots, energy statistics, and mesh-quality indicators hosted on the Exa-MA knowledge base (starting Q3~2026).
\end{itemize}